\subsection{内乘：代数情形}

设 \(E = \bigoplus_{p \geqslant 0} E_p\) 是一个 N 型分次 A-代数，P 是一个 Z 型分次 A-模；对每个齐次元素 \(x \in E_p\) ，用 \(x\) 左乘是 E 到自身的 A-线性映射 \(e(x)\) ，它是 \(p\) 次分次映射。对每个元素 \(u \in \text{Homgr}_A(E, P)\) ， \(u\) 被 \(x\) 的右内乘记作 \(u \llcorner x\) ，它就是元素 \(u \circ e(x)\) （属于 \(\text{Homgr}_A(E, P)\) ）。我们还记 \((i(x))(u) = u \llcorner x\) ，于是可见 \(i(x)\) 是 \(p\) 次分次自同态（分次 A-模 \(\text{Homgr}_A(E, P)\) 的）。若现在 \(x = \sum_{p \geqslant 0} x_p\) 是 E 的任意元素（其中 \(x_p \in E_p\) （对一切 \(p \geqslant 0\) ），且 \(x_p = 0\) （除 \(p\) 的有限个值外）），我们记 \(i(x) = \sum_{p=0}^{\infty} i(x_p)\) ，它因此是 A-模 \(\text{Homgr}_A(E, P)\) 的自同态。

为了记住在表达式 \(u \llcorner x\) 中哪个元素“作用”于另一个元素，注意“作用”于 u 的元素 x 位于 \(\llcorner\) 中水平线的自由端。

代数 E 的结合性转化为关系 \(e(xy) = e(x) \circ e(y)\) （对齐次的 x, y）；因此，根据 \(i(x)\) 的定义，

\[
i (x y) = i (y) \circ i (x)\tag{35}
\]

首先对齐次的 \(x, y\) ，然后通过线性性，对 E 中任意的 \(x, y\) ；这也可以写成

\[
(u \llcorner x) \llcorner y = u \llcorner (x y)\tag{36}
\]

对于 \(x, y\) （在 \(\mathbf{E}\) 中）及 \(u \in \operatorname{Homgr}_{\mathbf{A}}(\mathbf{E}, \mathbf{P})\) ；另一方面，显然 \(i(1)\) 是恒等映射（因为由 \(e(1) = 1_{\mathbf{E}}\) 即得），且 \(x \mapsto i(x)\) 是 A-线性的，于是可以看出，外合成律 \((x, u) \mapsto u \llcorner x\) （ \(x \in \mathbf{E}\) ， \(u \in \operatorname{Homgr}_{\mathbf{A}}(\mathbf{E}, \mathbf{P})\) ）与加法一起，在 \(\operatorname{Homgr}_{\mathbf{A}}(\mathbf{E}, \mathbf{P})\) 上定义了一个右 E-模结构。

特别地，我们考虑 P = A 的情形， \(\operatorname{Homgr}_{\mathbf{A}}(\mathbf{E}, \mathbf{P})\) 在此情形下是 E 的分次对偶 \(E^{*gr}\) ； \(i(x)\) 则是 A-线性映射 \(e(x)\) 的分次转置（II，§11，第6号），换言之，对 E 中所有 x, y 及 \(u \in E^{*gr}\) ，

\[
\langle u \llcorner x, y \rangle = \langle u, x y \rangle .\tag{37}
\]

按照本段开头的约定，注意，若 \(x \in E_{p}\) ，则 \(i(x)\) 是 \(E^{*gr}\) 的 -p 次自同态。

对每个齐次元素 \(x \in \mathbf{E}_p\) ，用 \(x\) 右乘类似地记作 \(e'(x)\) ，而元素 \(u \circ e'(x)\) （属于 \(\operatorname{Homgr}_{\mathbf{A}}(\mathbf{E}, \mathbf{P})\) ）称为 \(u\) 被 \(x\) 的左内乘，记作 \(x \lrcorner u\) ；我们记 \(i'(x) = x \lrcorner u\) ，于是 \(i'(x)\) 是 \(\operatorname{Homgr}_{\mathbf{A}}(\mathbf{E}, \mathbf{P})\) 的 \(p\) 次分次自同态；如上所述，此定义可推广到 \(x\) 是 \(\mathbf{E}\) 的任意元素的情形。由于这时 \(e'(xy) = e'(y)e'(x)\) ，

\[
i ^ {\prime} (x y) = i ^ {\prime} (x) \circ i ^ {\prime} (y)\tag{38}
\]

也可以写作

\[
x \lrcorner (y \lrcorner u) = (x y) \lrcorner u\tag{39}
\]

并且表明外合成律 \((x,u)\mapsto x\lrcorner u\) 与加法一起，在 \(\operatorname{Homgr}_{\mathbf{A}}(\operatorname{E},\operatorname{P})\) 上定义了一个左 E-模结构。另一方面，E 的结合性蕴含 \(e(x)\circ e'(y)=e'(y)\circ e(x)\) （对 E 中齐次的 x,y），由此得到关系式

\[
(y \lrcorner u) \llcorner x = y \lrcorner (u \llcorner x)\tag{40}
\]

从而使 \(\mathrm{Homgr}_{\mathbf{A}}(\mathbf{E},\mathbf{P})\) 上的两个外合成律在此集合上定义一个 (E, E)-双模结构（II，§ 1，第14号）。

当我们取 \(\mathbf{P} = \mathbf{A}\) 时， \(i'(x)\) 是 \(e'(x)\) 的分次转置；换言之，对所有 \(x, y\) （在 \(\mathbf{E}\) 中）及 \(u \in \mathbf{E}^{*\mathrm{gr}}\) ，

\[
\langle y, x \lrcorner u \rangle = \langle y x, u \rangle .\tag{41}
\]

当分次代数 E 是交换的时，显然 \(u \llcorner x = x \lrcorner u\) 。当 E 是反交换的且 P = A 时，则对于 \(x \in E_{p}\) 、 \(y \in E_{r}\) 和 \(u \in E_{q}^{*}\) ，有 \(yx = (-1)^{pr}xy\) ，由此，由 (37) 和 (41) 得 \(\langle u \llcorner x, y \rangle = (-1)^{pr}\langle y, x \lrcorner u \rangle\) 。但由于该关系的两边除 r = q - p 外均为零，

\[
x \lrcorner u = (- 1) ^ {p (q - p)} u \llcorner x.
\]

设 F 为另一个分次 A-代数，且 \(\phi: E \to F\) 为分次代数的一个 A-同态；则 \(\tilde{\phi} = \operatorname{Hom}(\phi, 1_{P}): \operatorname{Homgr}_{A}(F, P) \to \operatorname{Homgr}_{A}(E, P)\) 是 0 次的分次 A-同态；根据定义，对于 E 中的 \(x, y\) 及 \(u \in \operatorname{Homgr}_{A}(F, P)\) ，

\[
\begin{array}{r l} (\tilde {\phi} (u \llcorner \phi (x))) (y) & = (u \llcorner \phi (x)) (\phi (y)) \\ & = u (\phi (x) \phi (y)) = u (\phi (x y)) = (\tilde {\phi} (u)) (x y) = (\tilde {\phi} (u) \llcorner x) (y) \end{array}
\]

或等价地

\[
\tilde {\phi} (u \llcorner \phi (x)) = \tilde {\phi} (u) \llcorner x\tag{42}
\]

并且类似地

\[
\tilde {\phi} (\phi (x) \lrcorner u) = x \lrcorner \tilde {\phi} (u).\tag{43}
\]

换句话说，当 \(\operatorname{Homgr}_{\mathbf{A}}(\mathbf{F}, \mathbf{P})\) 借助环同态 \(\phi: \mathbf{E} \to \mathbf{F}\) 被视为 (E, E)-双模时，可以看出 \(\tilde{\phi}\) 是一个 (E, E)-双模同态（或者是把 (F, F)-双模 \(\operatorname{Homgr}_{\mathbf{A}}(\mathbf{F}, \mathbf{P})\) 映入 (E, E)-双模 \(\operatorname{Homgr}_{\mathbf{A}}(\mathbf{E}, \mathbf{P})\) 的 E-同态）。

例子。特别地，当 E 是分次代数 \(\mathsf{T}(\mathbf{M})\) 、 \(\mathsf{S}(\mathbf{M})\) 或 \(\bigwedge(\mathbf{M})\) 之一（其中 M 是一个 A-模），而 P 是一个 A-模（带有平凡分次）时，上述内容可以应用。为了显式求出如此得到的双模结构，注意 \(\operatorname{Homgr}_{\mathbf{A}}(\mathsf{T}(\mathbf{M}),\mathbf{P})\) （相应地 \(\operatorname{Homgr}_{\mathbf{A}}(\mathsf{S}(\mathbf{M}),\mathbf{P})\) 、 \(\operatorname{Homgr}_{\mathbf{A}}(\bigwedge(\mathbf{M}),\mathbf{P})\) ）中 -n 次的元素与从 \(M^{n}\) 到 P 的 n 重线性映射（相应地，对称 n 重线性映射、交错 n 重线性映射）等同。只需表出乘积

\[
f \llcorner (x _ {1} \otimes x _ {2} \otimes \dots \otimes x _ {p})
\]

\((\text{resp. } f \llcorner (x_1 x_2 \ldots x_p), \text{ resp. } f \llcorner (x_1 \wedge x_2 \wedge \dots \wedge x_p))\) （对 M 中元素的每个有限序列 \((x_i)_{1 \leq i \leq p}\) ）以及左内乘的类似表达式。由定义立即得出

\[
f \llcorner (x _ {1} \otimes x _ {2} \otimes \dots \otimes x _ {p}) = (x _ {1} \otimes x _ {2} \otimes \dots \otimes x _ {p}) \lrcorner f = 0\tag{44}
\]

若 \(p > n\) ，并且对于 \(p \leqslant n, f \llcorner (x_1 \otimes x_2 \otimes \dots \otimes x_p)\) （相应地

\[
\left(x _ {1} \otimes x _ {2} \otimes \dots \otimes x _ {p}\right) \lrcorner f)
\]

是由下式定义的 \((n - p)\) -线性映射：

\[
\left\{ \begin{array}{l} (f \llcorner (x _ {1} \otimes x _ {2} \otimes \dots \otimes x _ {p})) (y _ {1}, \ldots , y _ {n - p}) \\ = f (x _ {1}, \ldots , x _ {p}, y _ {1}, \ldots , y _ {n - p}) \\ ((x _ {1} \otimes x _ {2} \otimes \dots \otimes x _ {p}) \lrcorner f) (y _ {1}, \ldots , y _ {n - p}) \\ = f (y _ {1}, \ldots , y _ {n - p}, x _ {1}, \ldots , x _ {p}). \end{array} \right.\tag{45}
\]

对于 \(p > n\) ，在 \(\operatorname{Homgr}_{\mathbf{A}}(\mathsf{S}(\mathbf{M}), \mathbf{P})\) （相应地 \(\operatorname{Homgr}_{\mathbf{A}}(\bigwedge(\mathbf{M}), \mathbf{P})\) ）中也有公式 (44)，其中 \(x_1 \otimes x_2 \otimes \cdots \otimes x_p\) 替换为 \(x_1x_2 \ldots x_p\) （相应地 \(x_1 \wedge x_2 \wedge \cdots \wedge x_p\) ）。对于 \(p \leqslant n\) ，在 (45) 中进行相同的替换，就定义了对称 \((n - p)\) -线性映射 \(f \llcorner (x_1x_2 \ldots x_p)\) 和 \((x_1x_2 \ldots x_p) \lrcorner f\) （相应地，交错 \((n - p)\) -线性映射

\[
f \llcorner (x _ {1} \wedge x _ {2} \wedge \dots \wedge x _ {p}) \quad \text { and } \quad (x _ {1} \wedge x _ {2} \wedge \dots \wedge x _ {p}) \lrcorner f).
\]

当 \(n = p\) 时，上述乘积等于 \(\mathbf{M}\) 上等于 \(f(x_{1},\ldots ,x_{p})\) 的常值函数。

若 \(u: \mathbf{M} \to \mathbf{N}\) 是 A-模同态，则 \(\mathsf{T}(u): \mathsf{T}(\mathbf{M}) \to \mathsf{T}(\mathbf{N})\) 是分次 A-代数同态，于是由我们上面所见可知， \((\mathsf{T}(u))^{\sim}\) 是 \(\mathsf{T}(\mathbf{M})\) -同态，把 \((\mathsf{T}(\mathbf{N}), \mathsf{T}(\mathbf{N}))\) -双模 \(\operatorname{Homgr}_{\mathbf{A}}(\mathsf{T}(\mathbf{N}), \mathbf{P})\) 映入 \((\mathsf{T}(\mathbf{M}), \mathsf{T}(\mathbf{M}))\) -双模 \(\operatorname{Homgr}_{\mathbf{A}}(\mathsf{T}(\mathbf{M}), \mathbf{P})\) ，相对于环同态 \(\mathsf{T}(u)\) 。对 \((\mathbf{S}(u))^{\sim}\) 和 \((\bigwedge(u))^{\sim}\) 有类似的结果。
% semantic-fragments: legacy-input-seam
\relax{}
